%==============================================================================
%  ACADEMIC CV  —  U.S. PhD APPLICATIONS
%  Tailored for a fresh B.Sc. graduate (BUET, EEE) with active research,
%  manuscripts in preparation, and research-grade competition work.
%
%  COMPILE:  pdflatex cv.tex   (run twice so page totals resolve)
%  OVERLEAF: New Project -> Upload Project -> drop this file in. Nothing else needed.
%
%  ------------------------------------------------------------------------
%  READ THIS BEFORE YOU SUBMIT ANYTHING
%  ------------------------------------------------------------------------
%  Every FACT below is a blank: XX, [Month Year], [Title], 3.XX/4.00, etc.
%  The structure and the sentence patterns are real and ready to use; the
%  numbers are deliberately not. Fill in your own or delete the line.
%  Never send a CV that still contains an XX. See README.md for what belongs
%  in each section and which sections to delete.
%==============================================================================

\documentclass[11pt,letterpaper]{article}   % letterpaper: U.S. standard, not A4

\usepackage[T1]{fontenc}
\usepackage[utf8]{inputenc}
\usepackage{charter}            % Body font. Warm, readable, not overused.
\usepackage{microtype}          % Subtle spacing polish
\usepackage[
  top=0.8in, bottom=0.8in, left=0.85in, right=0.85in, footskip=0.35in
]{geometry}
\usepackage{titlesec}
\usepackage{enumitem}
\usepackage{xcolor}
\usepackage{tabularx}
\usepackage{fancyhdr}
\usepackage{lastpage}
\usepackage{needspace}          % Stops a section heading stranding at a page bottom
\usepackage[colorlinks=true,urlcolor=accent,linkcolor=accent]{hyperref}

%------------------------------------------------------------------------------
%  IDENTITY — edit these three lines and they propagate everywhere
%------------------------------------------------------------------------------
\newcommand{\myname}{Rifat Arafat}          % EXACTLY as printed in your passport
\newcommand{\myemail}{rifatarafat1020@gmail.com}
\newcommand{\myphone}{+880 1XXX-XXXXXX}

\definecolor{accent}{HTML}{1F3A5F}          % Deep navy. Restrained on purpose.
\definecolor{subtle}{HTML}{555555}

\hypersetup{pdftitle={\myname{} — Curriculum Vitae}, pdfauthor={\myname}}

%------------------------------------------------------------------------------
%  SECTION + ENTRY STYLING
%------------------------------------------------------------------------------
\titleformat{\section}
  {\needspace{4.5\baselineskip}\normalfont\normalsize\bfseries\scshape\color{accent}}
  {}{0em}{}
  [\vspace{-7pt}{\color{accent}\rule{\textwidth}{0.7pt}}]
\titlespacing*{\section}{0pt}{11pt}{5pt}

\setlist[itemize]{
  leftmargin=1.3em, itemsep=1.2pt, parsep=0pt, topsep=3pt,
  label=\textcolor{accent}{\small$\bullet$}
}

% \entry{Role or Degree}{Dates}{Organization}{Location}
\newcommand{\entry}[4]{%
  \vspace{3pt}%
  \noindent\begin{tabularx}{\textwidth}{@{}X r@{}}
    \textbf{#1} & \textbf{#2}\\
    \textit{\textcolor{subtle}{#3}} & \textit{\textcolor{subtle}{#4}}\\
  \end{tabularx}%
  \vspace{-4pt}%
}

% \line{Left text}{Right text} — one-line rows (awards, service, talks)
\newcommand{\lineentry}[2]{%
  \noindent\begin{tabularx}{\textwidth}{@{}X r@{}}
    #1 & #2\\
  \end{tabularx}%
}

\newcommand{\me}[1]{\underline{#1}}          % Underline your own name in citations
\newcommand{\skillrow}[2]{%
  \noindent\parbox[t]{0.21\textwidth}{\textbf{#1}}%
  \parbox[t]{0.78\textwidth}{#2}\par\vspace{2.5pt}%
}

\pagestyle{fancy}
\fancyhf{}
\renewcommand{\headrulewidth}{0pt}
\fancyfoot[C]{\footnotesize\textcolor{subtle}{\myname{} \quad$\cdot$\quad Curriculum Vitae \quad$\cdot$\quad Page \thepage{} of \pageref{LastPage}}}

\setlength{\parindent}{0pt}

%==============================================================================
\begin{document}
%==============================================================================

%------------------------------------------------------------------------------
%  HEADER
%  No photo. No date of birth. No marital status, religion, father's name, or
%  national ID. These are normal on a Bangladeshi CV and are a liability on a
%  U.S. academic one — committees read them as a signal you don't know the norms.
%------------------------------------------------------------------------------
\begin{center}
  {\LARGE\bfseries\color{accent}\myname}\\[5pt]
  {\small
    Dhaka, Bangladesh \quad$\cdot$\quad
    \href{mailto:\myemail}{\myemail} \quad$\cdot$\quad
    \myphone
  }\\[3pt]
  {\small
    % Delete any link you don't have. An empty Scholar profile is worse than none.
    \href{https://yourusername.github.io}{Website} \quad$\cdot$\quad
    \href{https://scholar.google.com/citations?user=XXXX}{Google Scholar} \quad$\cdot$\quad
    \href{https://github.com/yourusername}{GitHub} \quad$\cdot$\quad
    \href{https://linkedin.com/in/yourusername}{LinkedIn}
  }
\end{center}
\vspace{2pt}

%------------------------------------------------------------------------------
%  RESEARCH INTERESTS
%  The single most-read block on the page. A faculty member skims this to decide
%  whether you belong in their lab. Make it specific enough that a wrong-fit
%  professor self-selects out. Rewrite it per school (see README).
%------------------------------------------------------------------------------
\section{Research Interests}
[Primary area, e.g.\ analog/mixed-signal IC design] \,$\cdot$\, [Second area] \,$\cdot$\, [Third area]

\smallskip
\noindent Specifically: [one or two sentences naming the actual problem you want to work on
— the physics, the circuit, the algorithm. ``Machine learning'' is not a research
interest; ``learned digital pre-distortion for mmWave power amplifiers'' is.]

%------------------------------------------------------------------------------
%  EDUCATION
%------------------------------------------------------------------------------
\section{Education}

\entry{B.Sc.\ in Electrical and Electronic Engineering}{[Month Year] -- [Month Year]}
      {Bangladesh University of Engineering and Technology (BUET)}{Dhaka, Bangladesh}
\begin{itemize}
  \item \textbf{CGPA:} 3.XX\,/\,4.00 \quad$\cdot$\quad \textbf{Class Rank:} XX of XXX
        % Rank carries real weight from BUET — U.S. committees know the school is
        % highly selective but not how your GPA maps. Delete if not in the top ~15%.
  \item \textbf{Undergraduate Thesis:} ``[Full thesis title]''\\
        Supervisor: [Prof.\ Full Name], [Title], Dept.\ of EEE, BUET.
  \item \textbf{Selected Coursework:} Digital Signal Processing; Semiconductor Devices;
        VLSI Design; Random Signals and Noise; Control Systems; Electromagnetic Fields
        and Waves; Power Electronics; Communication Systems.
        % List only courses that support your stated research interest.
  \item \textbf{Relevant Graduate-Level / Self-Study:} [e.g.\ Convex Optimization
        (Stanford EE364A, audited); Analog IC Design (Razavi lecture series)].
        % Delete unless you genuinely completed it and can discuss it in an interview.
\end{itemize}

%------------------------------------------------------------------------------
%  RESEARCH EXPERIENCE
%  Your strongest section. With no publications yet, this is what convinces a
%  committee you can do research. Depth over breadth: three entries described
%  concretely beat eight one-liners.
%
%  Bullet formula:  ACTION VERB + technical specifics + method/tool + measured outcome
%  Weak:   "Worked on machine learning for signal processing."
%  Strong: "Built a 1-D CNN classifier for XX-channel EMG, reaching XX% accuracy
%           on XX subjects — XX points above the wavelet-feature baseline."
%------------------------------------------------------------------------------
\section{Research Experience}

\entry{Research Assistant}{[Month Year] -- Present}
      {[Lab / Group Name], Dept.\ of EEE, BUET \quad$\cdot$\quad Advisor: [Prof.\ Name]}{Dhaka, Bangladesh}
\begin{itemize}
  \item Designed / derived / implemented [the specific artifact], achieving [measured
        result] against [the baseline or ground truth you compared to].
  \item Second concrete contribution: [the hard technical problem you personally solved,
        not what the lab does in general].
  \item Validated results against [COMSOL / Cadence Spectre / measured hardware /
        published dataset]; [state the agreement, e.g.\ within X\% across N cases].
  \item \textit{Manuscript in preparation for [target venue]; expected submission [Month Year].}
        % Only if a draft genuinely exists. See README — this line is checkable.
\end{itemize}

\entry{Undergraduate Thesis Researcher}{[Month Year] -- [Month Year]}
      {Dept.\ of EEE, BUET \quad$\cdot$\quad Supervisor: [Prof.\ Name]}{Dhaka, Bangladesh}
\begin{itemize}
  \item Investigated [the research question] using [method], addressing [the specific
        gap in prior work — name the limitation you were attacking].
  \item Key result: [what you found, quantified].
  \item Presented findings to [thesis committee / departmental seminar] of [XX] faculty and students.
\end{itemize}

% Optional third entry: summer research, remote collaboration, or a professor's
% project you joined. Delete if you have none — a short honest section is fine.
\entry{[Research Intern / Project Collaborator]}{[Month Year] -- [Month Year]}
      {[Institution or Group] \quad$\cdot$\quad Advisor: [Name]}{[City, Country]}
\begin{itemize}
  \item Contribution: [what you did and what it produced, quantified].
\end{itemize}

%------------------------------------------------------------------------------
%  PUBLICATIONS AND MANUSCRIPTS
%
%  You have no published papers yet. That is completely normal for a U.S. PhD
%  applicant from a B.Sc. — publications are a bonus, not a requirement, and
%  engineering committees expect fewer of them than the life sciences do.
%
%  Keep this section ONLY if you have at least one genuine draft. If everything
%  is still an idea, delete the whole section and let the "manuscript in
%  preparation" line inside Research Experience carry it. An empty or padded
%  publication section reads worse than no section at all.
%------------------------------------------------------------------------------
\section{Manuscripts in Preparation}
{\small\textcolor{subtle}{Applicant's name underlined.}}
\begin{enumerate}[leftmargin=2.2em,label={[M\arabic*]},itemsep=3pt,topsep=4pt]
  \item \me{R.~Arafat}, [Co-author], and [Advisor]. ``[Working title].''
        \textit{In preparation for [IEEE Transactions on XXX / conference];
        targeted submission [Month Year].}
  \item \me{R.~Arafat} and [Advisor]. ``[Working title].'' \textit{In preparation.}
\end{enumerate}

% When a paper is actually submitted, move it into a "Under Review" list and
% cite it properly. When it is accepted, it graduates to "Publications" at the
% very top of this section. Never list an idea as if it were a submission.

%------------------------------------------------------------------------------
%  RESEARCH COMPETITIONS AND TECHNICAL PROJECTS
%  This is where your competition work belongs — and it is a real asset, because
%  it shows independent problem-selection under time pressure. Describe it the
%  way you would describe research: what was novel, what you measured, where you
%  placed. Placement matters; so does the size of the field you beat.
%------------------------------------------------------------------------------
\section{Research Competitions and Technical Projects}

\entry{[Competition Name] --- [Placement, e.g.\ 1st of XXX teams]}{[Month Year]}
      {[Organizer, e.g.\ IEEE Region 10 / national round]}{[Location]}
\begin{itemize}
  \item \textbf{Novel contribution:} [State plainly what had not been done before —
        the new topology, the new formulation, the constraint you removed.]
  \item Built [the system] using [stack/tools]; demonstrated [measured performance]
        under [the competition's constraints].
  \item Led a team of [X]; personally owned [your specific subsystem].
\end{itemize}

\entry{[Project Title]}{[Month Year]}
      {[Course project / independent work] \quad$\cdot$\quad \href{https://github.com/yourusername/repo}{Code}}{}
\begin{itemize}
  \item What it does, what was technically hard, and the result: [quantified].
\end{itemize}

%------------------------------------------------------------------------------
%  HONORS AND AWARDS
%------------------------------------------------------------------------------
\section{Honors and Awards}
\lineentry{[University Merit Scholarship / Dean's List], BUET \hfill}{[Year]}
\lineentry{[Award name], [awarding body] --- [selectivity, e.g.\ XX of XXX recipients nationally] \hfill}{[Year]}
\lineentry{[Competition award, if not already listed above] \hfill}{[Year]}
\lineentry{[Board Scholarship, HSC/SSC — Talentpool or General] \hfill}{[Year]}
% Keep national-level and university-level honors. Drop school club certificates.

%------------------------------------------------------------------------------
%  TECHNICAL SKILLS
%  Tools you could be handed in an interview and asked to use. Nothing aspirational.
%  Skip self-ratings ("Python: 8/10") — U.S. committees find them odd.
%------------------------------------------------------------------------------
\section{Technical Skills}
\skillrow{Programming}{Python, MATLAB, C/C++, Verilog/SystemVerilog, [VHDL], [Bash]}
\skillrow{Frameworks}{PyTorch, NumPy/SciPy, scikit-learn, [Simulink]}
\skillrow{EDA \& Simulation}{Cadence Virtuoso/Spectre, LTspice, [Synopsys], [COMSOL], [HFSS/ADS]}
\skillrow{Hardware}{[FPGA: Xilinx Vivado], [STM32/Arduino], oscilloscope and VNA measurement, PCB design [Altium/KiCad]}
\skillrow{Other}{\LaTeX, Git, Linux, [Origin/Excel]}

%------------------------------------------------------------------------------
%  TEACHING AND MENTORING
%  Keep this if you have it — U.S. PhD funding is usually a TA-ship, so evidence
%  that you can teach in English is directly relevant. Delete if you have none.
%------------------------------------------------------------------------------
\section{Teaching and Mentoring}
\entry{[Undergraduate Teaching Assistant / Tutor]}{[Month Year] -- [Month Year]}
      {[Course name and code], Dept.\ of EEE, BUET}{Dhaka, Bangladesh}
\begin{itemize}
  \item Led [weekly lab sessions / tutorials] for [XX] students; graded [assignments/labs].
  \item Mentored [X] junior students on [topic]; [outcome, if any].
\end{itemize}

%------------------------------------------------------------------------------
%  TEST SCORES
%  Include only what is strong. Many U.S. programs have dropped the GRE entirely
%  — check each one. TOEFL/IELTS is still required of most international
%  applicants. A weak score is better left off the CV (the portal gets it anyway).
%------------------------------------------------------------------------------
\section{Test Scores}
\lineentry{\textbf{IELTS Academic:} X.X overall (L X.X, R X.X, W X.X, S X.X) \hfill}{[Month Year]}
\lineentry{\textbf{TOEFL iBT:} XXX (R XX, L XX, S XX, W XX) \hfill}{[Month Year]}
\lineentry{\textbf{GRE General:} Q XXX (XX\%), V XXX (XX\%), AW X.X \hfill}{[Month Year]}
% Delete the rows you did not take. Do not list a planned test date.

%------------------------------------------------------------------------------
%  PROFESSIONAL SERVICE AND MEMBERSHIPS  (optional — delete if thin)
%------------------------------------------------------------------------------
\section{Professional Service and Memberships}
\lineentry{Student Member, IEEE [and IEEE XXX Society] \hfill}{[Year] -- Present}
\lineentry{[Role], [IEEE BUET Student Branch / departmental society] \hfill}{[Year] -- [Year]}
\lineentry{[Reviewer / volunteer role for a conference, if any] \hfill}{[Year]}

%------------------------------------------------------------------------------
%  REFERENCES
%  Three is standard. Ask permission first, and give each of them this CV plus
%  your statement of purpose. Your thesis supervisor should be first.
%------------------------------------------------------------------------------
\section{References}
\vspace{2pt}
\noindent
\begin{tabularx}{\textwidth}{@{}XX@{}}
\textbf{[Prof.\ Full Name]} & \textbf{[Prof.\ Full Name]}\\
{[Professor]}, Dept.\ of EEE & {[Associate Professor]}, Dept.\ of EEE\\
Bangladesh University of Engineering \& Technology & Bangladesh University of Engineering \& Technology\\
\href{mailto:email@eee.buet.ac.bd}{email@eee.buet.ac.bd} & \href{mailto:email@eee.buet.ac.bd}{email@eee.buet.ac.bd}\\
\textit{Undergraduate thesis supervisor} & \textit{Course instructor and project advisor}\\[10pt]
\textbf{[Prof.\ Full Name]} & \\
{[Title]}, {[Department]} & \\
{[Institution]} & \\
\href{mailto:email@example.edu}{email@example.edu} & \\
\textit{[Relationship to you]} & \\
\end{tabularx}

\end{document}
